\section{Introduzione}
\subsection{Scopo del documento}
Il presente documento ha lo scopo di definire norme, convenzioni e strumenti che ogni membro del gruppo \textbf{\nomegruppo} seguirà per tutta la durata del progetto. La documentazione sarà aggiornata costantemente man mano che il gruppo adotta nuove pratiche. 
Le norme di progetto hanno i seguenti obiettivi:
\begin{itemize}
	\item Definire il way of working del gruppo;
	\item Garantire l'uniformità della documentazione
	\item Favorire il miglioramento continuo
	\item Standardizzare le attività
\end{itemize}

\subsection{Convenzioni Tipografiche}
Si adottano le seguenti convenzioni, che i membri del gruppo dovranno seguire in maniera rigorosa durante la stesura dei documenti:
\begin{itemize}
	\item I termini definiti nel glossario sono marcati con il pedice G;
	\item Riferimenti a file, branch, comandi dovranno essere riportati in font \texttt{monospace}.
	\item Riferimenti interni alla documentazione devono essere riportati attraverso l'utilizzo di hyperlink e usando il font \textit{corsivo}.
\end{itemize}

\subsection{Riferimenti}
\subsubsection{Normativi}
\begin{itemize}
		\item Regolamento del Progetto - Corso di Ingegneria del Software A.A. 2026/2027
\end{itemize}

\subsubsection{Informativi}
\begin{itemize}
	\item \href{https://www.iso.org/iso-8601-date-and-time-format.html}{ISO 8601, \textit{Date and time}}
	\item \href{https://semver.org/}{T. Preston-Werner, \textit{Semantic Versioning 2.0.0}}: adottato come base per il formato \texttt{X.Y.Z};
	\item \href{https://www.conventionalcommits.org/en/v1.0.0/}{Conventional Commits 1.0.0}
\end{itemize}
